\section{Design Decisions}

\subsection{Language Overview}

The language is an imperative, Rust-inspired language whose central novelty is
an \emph{ownership type system}: every value has exactly one owner at any time,
non-copyable values cannot be used after being moved, and aliasing is controlled
through a borrow system with static exclusivity guarantees.
The full feature set covers ten phases:

\begin{itemize}
  \item \textbf{Phase 0}: mutable/immutable variable bindings, block scoping,
    \texttt{if}/\texttt{while} statements, multi-parameter functions
  \item \textbf{Phase 1}: affine \texttt{Light} type (traffic-light enum) with
    move semantics --- values are consumed on read and cannot be used again
  \item \textbf{Phase 2A/B/C}: lists, copyable primitives (\texttt{int},
    \texttt{bool}), \texttt{Result}, pairs, and pattern matching
  \item \textbf{Phase 3A/B}: immutable borrows (\texttt{\&x}), mutable borrows
    (\texttt{\&mut x}) with exclusivity (at most one mutable borrow, or any
    number of immutable borrows, but never both simultaneously)
  \item \textbf{Phase 3C}: non-lexical lifetimes (NLL) --- borrows expire at
    the \emph{last syntactic use}, not at the enclosing block boundary
  \item \textbf{Phase 4A}: explicit lifetime annotations on functions that
    return references (\texttt{fn f<'a>(x: \&'a T) -> \&'a T})
  \item \textbf{Phase 4B}: type-safe \texttt{spawn} blocks --- the type checker
    ensures only \texttt{Copy} values are captured from the enclosing scope
\end{itemize}

\subsection{Architecture}

The implementation is split into a grammar package (BNFC-generated lexer and
LALR(1) parser) and a library package containing the type checker and
interpreter.
Both the checker and the interpreter share the same \texttt{ScopeStack}
data structure and use the same lens-based context-update idiom.

\begin{figure}[ht!]
    \centering
    \begin{minipage}{.9\linewidth}
\begin{lstlisting}
grammar/Lang.cf        BNFC grammar (99 rules); Alex lexer + Happy parser
src/ScopeStack.hs      Persistent scope stack
src/Value.hs           Runtime values; TClosure; isCopyable; eraseLifetime
src/Context.hs         TcCtx, EvalCtx records; van Laarhoven lenses
src/Tc.hs              Tc  = ExceptT String (State TcCtx)
src/Eval.hs            Eval = ExceptT String (State EvalCtx)
src/Logger.hs          Background logging thread (forkIO + Chan + MVar)
src/TypeCheck/Stmt.hs  Statement checking: borrows, NLL, lifetimes, spawn
src/TypeCheck/Expr.hs  Expression type inference: ownership, EMatch
src/Interp/Stmt.hs     Statement evaluation
src/Interp/Expr.hs     Expression evaluation: ECall, EDeref, EMatch
\end{lstlisting}
    \end{minipage}
    \captionof{lstlisting}{Module layout}
    \label{lst:modules}
\end{figure}

\subsection{Per-Variable Ownership Tracking}

The type checker associates each in-scope variable with a \texttt{VarInfo}
record that encodes its full ownership state:

\begin{figure}[ht!]
    \centering
    \begin{minipage}{.9\linewidth}
\begin{lstlisting}[language=Haskell]
data VarInfo = VarInfo
  { varType       :: Type        -- declared type
  , varMut        :: Bool        -- declared with mut?
  , varOwned      :: Bool        -- False after a move
  , varBorrows    :: Int         -- active immutable borrow count
  , varMutBorrows :: Int         -- active mutable borrow count (max 1)
  , varBorrowOf   :: Maybe Ident -- referent if this var is a borrow
  }
\end{lstlisting}
    \end{minipage}
    \captionof{lstlisting}{\texttt{VarInfo} --- per-variable state in the type checker}
    \label{lst:varinfo}
\end{figure}

When a non-\texttt{Copy} variable is read (e.g.\ passed to a function or bound
via \texttt{let}), \texttt{varOwned} is set to \texttt{False}.
Any subsequent read triggers a ``used after being moved'' error.
The \texttt{varBorrows} and \texttt{varMutBorrows} counters enforce exclusivity:
a mutable borrow requires both counters to be zero; an immutable borrow requires
only the mutable counter to be zero.

\subsection{Example Programs}

\Cref{lst:example} shows a program combining mutable borrows and NLL.
The immutable borrow \texttt{borrow} expires at its last syntactic use
(\texttt{*borrow}), which allows the subsequent mutable borrow inside
\texttt{set\_red} to proceed without error.

\begin{figure}[ht!]
    \centering
    \begin{minipage}{.9\linewidth}
        \lstinputlisting{assets/examples/sample.afp}
    \end{minipage}
    \captionof{lstlisting}{Mutable borrows and NLL, returning \texttt{Red}}
    \label{lst:example}
\end{figure}

\Cref{lst:lt-example} shows explicit lifetime annotations (Phase 4A).
The type checker resolves the return borrow at the call site: \texttt{r}
borrows \texttt{a} (via lifetime \texttt{'a}), so \texttt{a} cannot be moved
while \texttt{r} is live.

\begin{figure}[ht!]
    \centering
    \begin{minipage}{.9\linewidth}
\begin{lstlisting}
fn first<'a,'b>(x: &'a int, y: &'b int) -> &'a int { x };
let a = 3;
let b = 4;
let r = first(&a, &b);  // r borrows a via lifetime 'a
*r
// -> 3
\end{lstlisting}
    \end{minipage}
    \captionof{lstlisting}{Explicit lifetime annotations, returning \texttt{3}}
    \label{lst:lt-example}
\end{figure}

\Cref{lst:spawn-example} shows \texttt{spawn} (Phase 4B).
The block captures \texttt{count} (type \texttt{int}, which is \texttt{Copy})
safely.
Replacing \texttt{count} with a \texttt{Light} variable would be rejected at
the \texttt{spawn} site.

\begin{figure}[ht!]
    \centering
    \begin{minipage}{.9\linewidth}
\begin{lstlisting}
let count = 5;
spawn { let local = count };   // OK: int is Copy
// let x = Red;
// spawn { let y = x }         // REJECTED: Light is not Copy
count
// -> 5
\end{lstlisting}
    \end{minipage}
    \captionof{lstlisting}{Type-safe spawn: only \texttt{Copy} captures allowed}
    \label{lst:spawn-example}
\end{figure}

\subsection{Key Design Decisions and Trade-offs}

\paragraph{ScopeStack over a flat environment.}
The scope is a persistent list of \texttt{Map Ident a} frames.
The main alternative --- a single flat map with depth tags --- would require
scanning the entire map to remove variables on block exit (O($n$) per exit)
vs.\ O(1) list tail for \texttt{pop}.
More importantly, \texttt{updateStack} propagates mutable assignments outward
by rewriting the \emph{first frame} that contains the key, which correctly
implements mutation of outer variables from inner blocks without breaking the
persistent structure.
The trade-off is that lookups traverse all frames
(O($d \cdot \log n$) for nesting depth $d$), but in practice $d$ is small
and the code is significantly cleaner than the flat-map alternative.

\paragraph{VarInfo counters vs.\ an explicit borrow set.}
An alternative design would maintain a separate set of
$(\text{borrower}, \text{borrowee})$ pairs.
We chose integer counters (\texttt{varBorrows}, \texttt{varMutBorrows}) on
the \emph{borrowee} side: they are simpler to check (exclusivity is
\texttt{varMutBorrows > 0}) and update atomically with \texttt{over tcVars}.
The \texttt{varBorrowOf} field on the borrower side compensates:
\texttt{releaseTopBorrows} uses it to find and decrement the correct counter
when a reference leaves scope.
The weakness is that the counter does not identify \emph{which} specific
reference holds a borrow, making error messages slightly less precise than
a set-based design would allow.

\paragraph{NLL: syntactic over-approximation vs.\ precise dataflow.}
Rust computes NLL via full liveness analysis over a control-flow graph.
Our implementation uses a conservative \emph{syntactic free-variable}
approximation: \texttt{mentionedVars} collects all identifiers that appear
textually in the remaining statements, including those inside loop bodies
and both branches of conditionals.
This means a borrow can stay live longer than strictly necessary when a
variable is mentioned in a dead branch, so the checker accepts \emph{fewer}
programs than a full NLL analysis would.
The trade-off is \emph{soundness} (no borrow-dangling errors are missed) at
the cost of completeness.
For a course project the approximation is sufficient and avoids constructing
a control-flow graph entirely.

\paragraph{Synchronous \texttt{spawn}.}
The interpreter runs \texttt{spawn} blocks synchronously by calling
\texttt{interpBlock} --- the same function used for ordinary block statements.
True concurrent execution would require lifting the \texttt{Eval} monad into
\texttt{IO} (i.e., using \texttt{ExceptT String (StateT EvalCtx IO)}), which
would break the pure, deterministic test suite.
The design separates concerns: the \emph{type-safety guarantee} (Copy-capture
enforcement) is completely independent of the scheduling strategy.
A real concurrent runtime could be substituted without changing the type
checker at all.

\paragraph{Hand-rolled lenses vs.\ the \texttt{lens} library.}
Five van Laarhoven lenses suffice for this project.
Writing them by hand avoids a large transitive dependency and Template Haskell
splice compilation, and makes the representation fully transparent.
The key payoff appears in \texttt{releaseExpiredBorrows}, where the same
\texttt{Map.traverseWithKey} that can \emph{map} over entries also acts as a
zero-allocation \emph{structural fold} when instantiated with the \texttt{Const}
applicative --- collecting all expired borrows in a single pass over the
top scope frame.
